% part: intuitionistic-logic
% chapter: soundness-completeness
% section: soundness-nd

\documentclass[../../../include/open-logic-section]{subfiles}

\begin{document}

\olfileid{int}{sc}{snd}

\olsection{સ્વાભાવિક નિગમનની યથાર્થતા}

હવે આપણે સંબંધાત્મક અર્થવિજ્ઞાનની દૃષ્ટિએ સ્વાભાવિક નિગમનની
યથાર્થતા સાબિત કરીશું. એટલે કે, ધારણાઓના ગણમાંથી
!!a{formula} !!{derivable} હોય, તો તે ધારણાઓનો ગણ તે
!!{formula}ને ફલિત કરે છે એમ બતાવીશું.

\begin{thm}[યથાર્થતા]
  \ollabel{thm:soundness}
  જો $\Gamma \Proves !A$ હોય, તો $ \Gamma \Entails !A$.
\end{thm}

\begin{proof}
  જો $\Gamma \Proves !A$ હોય, તો $\Gamma \Entails !A$
  તે સાબિત કરીએ છીએ. $\Gamma$માંથી~$!A$ના
  !!{derivation} પર આગમન કરીએ છીએ.

  \begin{enumerate}
  \item જો !!{derivation}માં માત્ર ધારણા~$!A$ જ હોય,
    તો $!A \Proves !A$ છે અને $!A \Entails !A$
    બતાવવાનું છે. ધારો કે $\mSat{M}{!A}[w]$.
    તો સ્પષ્ટ છે કે $\mSat{M}{!A}[w]$.

  \item !!{derivation}નો છેલ્લો નિયમ $\Intro{\land}$ છે:
    \iftag{probAnd}{સ્વાધ્યાય.}{!!{undischarged} ધારણાઓ~$\Gamma$માંથી
      $!B$ની અને !!{undischarged} ધારણાઓ~$\Delta$માંથી~$!C$ની
      પૂર્વધારણાઓના !!{derivation}s,
      $\Gamma \Proves !B$ અને $\Delta \Proves !C$ બતાવે છે.
      આગમનની ધારણાથી $\Gamma \Entails !B$ અને
      $\Delta \Entails !C$ છે. આખા !!{derivation}ની
      !!{undischarged} ધારણાઓ $\Gamma$ અને $\Delta$ બંને
      છે, તેથી $\Gamma \cup \Delta \Entails !B \land !C$
      બતાવવાનું છે.
      \sourcecorrection{OLMD-090}{મૂળમાં અહીં નિષ્કર્ષ
      $!A \land !B$ લખાયો છે, જ્યારે પૂર્વધારણાઓ
      $!B$ અને $!C$ છે અને અંતે $!B \land !C$ મળે છે.}
      ધારો કે $\mSat{M}{\Gamma \cup \Delta}[w]$.
      ત્યારે $\mSat{M}{\Gamma}[w]$ પણ છે.
      $\Gamma \Entails !B$ હોવાથી $\mSat{M}{!B}[w]$.
      એ જ રીતે $\mSat{M}{!C}[w]$. તેથી
      $\mSat{M}{!B \land !C}[w]$.}

  \item !!{derivation}નો છેલ્લો નિયમ $\Elim{\land}$ છે:
    \iftag{probAnd}{સ્વાધ્યાય.}{પૂર્વધારણા
      $!B \land !C$નું !!{derivation}, જેમાં $\Gamma$ની
      ધારણાઓ !!{undischarged} છે, $\Gamma \Proves !B \land !C$
      બતાવે છે. આગમનની ધારણાથી $\Gamma \Entails !B \land !C$.
      $\Gamma \Entails !B$ બતાવવાનું છે. તેથી ધારો કે
      $\mSat{M}{\Gamma}[w]$. $\Gamma \Entails !B \land !C$
      હોવાથી $\mSat{M}{!B \land !C}[w]$; તેથી
      $\mSat{M}{!B}[w]$ પણ છે. એ જ રીતે જો $\Elim\land$નો
      નિષ્કર્ષ~$!C$ હોય, તો $\Gamma \Entails !C$.}

  \item !!{derivation}નો છેલ્લો નિયમ $\Intro{\lor}$ છે:
    \iftag{probOr}{સ્વાધ્યાય.}{ધારો કે પૂર્વધારણા $!B$
      છે અને $!B$ પર પૂરાં થતા !!{derivation}ની
      !!{undischarged} ધારણાઓ $\Gamma$ છે. ત્યારે
      $\Gamma \Proves !B$ અને આગમનની ધારણાથી
      $\Gamma \Entails !B$. $\Gamma \Entails !B \lor !C$
      બતાવવાનું છે. ધારો કે $\mSat{M}{\Gamma}[w]$.
      $\Gamma \Entails !B$ હોવાથી $\mSat{M}{!B}[w]$;
      તેથી $\mSat{M}{!B \lor !C}[w]$ પણ છે. એ જ રીતે
      પૂર્વધારણા~$!C$ હોય, તો $\Gamma \Entails !C$.}

  \item !!{derivation}નો છેલ્લો નિયમ~$\Elim{\lor}$ છે:
    \iftag{probOr}{સ્વાધ્યાય.}{પૂર્વધારણાઓ પર પૂરાં થતાં
      !!{derivation}s અનુક્રમે !!{undischarged} ધારણાઓ~$\Gamma$
      પરથી $!B \lor !C$નાં, !!{undischarged} ધારણાઓ
      $\Delta_1 \cup \{!B\}$ પરથી $!D$નાં અને
      !!{undischarged} ધારણાઓ~$\Delta_2 \cup \{!C\}$
      પરથી $!D$નાં છે. તેથી $\Gamma \Proves !B \lor !C$,
      $\Delta_1 \cup \{!B\} \Proves !D$, અને
      $\Delta_2 \cup \{!C\} \Proves !D$.
      આગમનની ધારણાથી $\Gamma \Entails !B \lor !C$,
      $\Delta_1 \cup \{!B\} \Entails !D$, અને
      $\Delta_2 \cup \{!C\} \Entails !D$.
      $\Gamma \cup \Delta_1 \cup \Delta_2 \Entails !D$
      સાબિત કરવાનું છે.

    ધારો કે $\mSat{M}{\Gamma \cup \Delta_1 \cup \Delta_2}[w]$.
    ત્યારે $\mSat{M}{\Gamma}[w]$ છે અને
    $\Gamma \Entails !B \lor !C$ હોવાથી
    $\mSat{M}{!B \lor !C}[w]$. $\mSat{M}{}$ની
    વ્યાખ્યા મુજબ $\mSat{M}{!B}[w]$ અથવા
    $\mSat{M}{!C}[w]$. તેથી બે કિસ્સા છે:
    (a)~$\mSat{M}{!B}[w]$. ત્યારે
    $\mSat{M}{\Delta_1 \cup \{!B\}}[w]$.
    \sourcecorrection{OLMD-091}{મૂળમાં (a)માં
    $\mSat{M}{!B}$ પછીનું $[w]$ ગણિત-મર્યાદાની
    બહાર છે; તેને એ જ સત્ય-સંજ્ઞામાં મૂક્યું છે.}
    $\Delta_1 \cup \{!B\} \Entails !D$ હોવાથી
    $\mSat{M}{!D}[w]$. (b)~$\mSat{M}{!C}[w]$.
    ત્યારે $\mSat{M}{\Delta_2 \cup \{!C\}}[w]$.
    $\Delta_2 \cup \{!C\} \Entails !D$ હોવાથી
    $\mSat{M}{!D}[w]$. આથી બંને કિસ્સામાં ઇચ્છિત
    $\mSat{M}{!D}[w]$ મળે છે.
    \sourcecorrection{OLMD-092}{મૂળના બે અનુસરણ-દાવામાં
    $!B$ અને $!C$ પહેલાં $\cup$ આવે છે, પરંતુ
    એકલ સૂત્રો ગણો નથી; અગાઉની પૂર્વધારણાઓ પ્રમાણે
    તેમને $\{!B\}$ અને $\{!C\}$ એકક ગણો કર્યા છે.}}

  \item !!{derivation}નો છેલ્લો નિયમ $\Intro{\lif}$ છે અને
    નિષ્કર્ષ~$!B\lif !C$ છે. ત્યારે પૂર્વધારણા~$!C$
    છે અને તેના પર પૂરાં થતા !!{derivation}ની
    !!{undischarged} ધારણાઓ~$\Gamma \cup \{!B\}$ છે.
    તેથી $\Gamma \cup \{!B\} \Proves !C$ અને આગમનની
    ધારણાથી $\Gamma \cup \{!B\} \Entails !C$.
    $\Gamma \Entails !B \lif !C$ બતાવવાનું છે.

    ધારો કે $\mSat{M}{\Gamma}[w]$. દરેક એવા $w'$ માટે
    કે $Rww'$, જો $\mSat{M}{!B}[w']$ હોય, તો
    $\mSat{M}{!C}[w']$ છે તે બતાવવું છે. તેથી
    $Rww'$ અને $\mSat{M}{!B}[w']$ ધારો.
    \olref[sem][rel]{prop:true-monotonic} મુજબ
    $\mSat{M}{\Gamma}[w']$. હવે
    $\Gamma \cup \{!B\} \Entails !C$ હોવાથી
    $\mSat{M}{!C}[w']$ મળે છે, જે બતાવવાનું હતું.

  \item !!{derivation}નો છેલ્લો નિયમ $\Elim{\lif}$ છે અને
    નિષ્કર્ષ~$!C$ છે. પૂર્વધારણાઓ $!B \lif !C$ અને
    $!B$ છે, જેમના !!{derivation}s અનુક્રમે
    !!{undischarged} ધારણાઓ $\Gamma$ અને $\Delta$
    પરથી છે. તેથી $\Gamma \Proves !B \lif !C$ અને
    $\Delta \Proves !B$. આગમનની ધારણાથી
    $\Gamma \Entails !B \lif !C$ અને
    $\Delta \Entails !B$. $\Gamma \cup \Delta
    \Entails !C$ બતાવવાનું છે.

    ધારો કે $\mSat{M}{\Gamma \cup \Delta}[w]$.
    $\mSat{M}{\Gamma}[w]$ અને
    $\Gamma \Entails !B \lif !C$ હોવાથી
    $\mSat{M}{!B \lif !C}[w]$. વ્યાખ્યા મુજબ,
    $Rww'$ ધરાવતા દરેક $w'$ માટે જો
    $\mSat{M}{!B}[w']$ હોય, તો
    $\mSat{M}{!C}[w']$. $R$ સ્વવાચક હોવાથી
    $Rww'$ ધરાવતા $w'$માં $w$ પોતે પણ આવે છે;
    એટલે જો $\mSat{M}{!B}[w]$ હોય, તો
    $\mSat{M}{!C}[w]$. હવે $\mSat{M}{\Delta}[w]$
    અને $\Delta \Entails !B$ હોવાથી
    $\mSat{M}{!B}[w]$. તેથી ઇચ્છિત
    $\mSat{M}{!C}[w]$ મળે છે.

  \item !!{derivation}નો છેલ્લો નિયમ $\FalseInt$ છે અને
    નિષ્કર્ષ~$!A$ છે. પૂર્વધારણા $\lfalse$ છે અને તેના
    !!{derivation}ની !!{undischarged} ધારણાઓ~$\Gamma$ છે.
    ત્યારે $\Gamma \Proves \lfalse$. આગમનની ધારણાથી
    $\Gamma \Entails \lfalse$. હવે $\Gamma \Entails !A$
    બતાવવાનું છે.

    પરોક્ષ રીતે આગળ વધીએ. જો $\Gamma \Entails/ !A$
    હોય, તો એવું નિદર્શ~$\mModel{M}$ અને જગત~$w$
    છે કે $\mSat{M}{\Gamma}[w]$ અને
    $\mSat/{M}{!A}[w]$. $\Gamma \Entails \lfalse$
    હોવાથી $\mSat{M}{\lfalse}[w]$. પરંતુ આ અશક્ય
    છે, કારણ કે વ્યાખ્યા મુજબ $\mSat/{M}{\lfalse}[w]$.
    તેથી $\Gamma \Entails !A$.
  \item !!{derivation}નો છેલ્લો નિયમ $\Intro\lnot$ છે: સ્વાધ્યાય.
  \item !!{derivation}નો છેલ્લો નિયમ $\Elim\lnot$ છે: સ્વાધ્યાય.
  \end{enumerate}
\end{proof}

\begin{prob}
  \olref[int][sc][snd]{thm:soundness}ની સાબિતી પૂરી કરો.
  $\Intro\lnot$ અને $\Elim\lnot$ના કિસ્સાઓમાં
  \olref[int][sem][rel]{defn:true-at-w}માં આપેલી
  $\mSat{M}{\lnot !A}[w]$ની વ્યાખ્યા વાપરો;
  એટલે $\lnot !A$ને $!A \lif \lfalse$ વડે
  વ્યાખ્યાયિત થયેલું માનીને કામ ન કરો.
\end{prob}

\begin{prob}
  નીચેના !!{formula}s સ્ફુરણાવાદી તર્કમાં !!{derivable}
  નથી તે બતાવો:
  \begin{enumerate}
    \item $(!A \lif !B) \lor (!B \lif !A)$
    \item $(\lnot\lnot !A \lif !A) \lif (!A \lor \lnot !A)$
    \item $(!A \lif !B \lor !C) \lif \bigl((!A \lif !B)\lor(!A \lif !C)\bigr)$
  \end{enumerate}
\end{prob}

\end{document}
